\documentclass[10pt,a4paper]{amsart}
\usepackage{amsmath,amssymb,amsthm,graphicx}
% Retained minimal preamble for the Stacks Project chapter sites-cohomology.tex (Cohomology on Sites).
% Verbatim excerpt of the chapter preamble: Xy-pic setup (native preamble lines 14-18) and the native
% theorem declarations with the mathematical macros (native preamble lines 154-217). Copied mechanically;
% nothing added, removed or redefined. The wrapper class and the amsmath/amssymb/amsthm/graphicx base
% packages are supplied by the assembler, so the class/externaldocument part of the native preamble is omitted.
% For commutative diagrams we use Xy-pic
\usepackage[all]{xy}

% We use 2cell for 2-commutative diagrams.
\xyoption{2cell}
% Theorem environments.
%
\theoremstyle{plain}
\newtheorem{theorem}[subsection]{Theorem}
\newtheorem{proposition}[subsection]{Proposition}
\newtheorem{lemma}[subsection]{Lemma}

\theoremstyle{definition}
\newtheorem{definition}[subsection]{Definition}
\newtheorem{example}[subsection]{Example}
\newtheorem{exercise}[subsection]{Exercise}
\newtheorem{situation}[subsection]{Situation}

\theoremstyle{remark}
\newtheorem{remark}[subsection]{Remark}
\newtheorem{remarks}[subsection]{Remarks}

\numberwithin{equation}{subsection}

% Macros
%
\def\lim{\mathop{\mathrm{lim}}\nolimits}
\def\colim{\mathop{\mathrm{colim}}\nolimits}
\def\Spec{\mathop{\mathrm{Spec}}}
\def\Hom{\mathop{\mathrm{Hom}}\nolimits}
\def\Ext{\mathop{\mathrm{Ext}}\nolimits}
\def\SheafHom{\mathop{\mathcal{H}\!\mathit{om}}\nolimits}
\def\SheafExt{\mathop{\mathcal{E}\!\mathit{xt}}\nolimits}
\def\Sch{\mathit{Sch}}
\def\Mor{\mathop{\mathrm{Mor}}\nolimits}
\def\Ob{\mathop{\mathrm{Ob}}\nolimits}
\def\Sh{\mathop{\mathit{Sh}}\nolimits}
\def\NL{\mathop{N\!L}\nolimits}
\def\CH{\mathop{\mathrm{CH}}\nolimits}
\def\proetale{{pro\text{-}\acute{e}tale}}
\def\etale{{\acute{e}tale}}
\def\QCoh{\mathit{QCoh}}
\def\Ker{\mathop{\mathrm{Ker}}}
\def\Im{\mathop{\mathrm{Im}}}
\def\Coker{\mathop{\mathrm{Coker}}}
\def\Coim{\mathop{\mathrm{Coim}}}

% Boxtimes
%
\DeclareMathSymbol{\boxtimes}{\mathbin}{AMSa}{"02}

%
% Macros for moduli stacks/spaces
%
\def\QCohstack{\mathcal{QC}\!\mathit{oh}}
\def\Cohstack{\mathcal{C}\!\mathit{oh}}
\def\Spacesstack{\mathcal{S}\!\mathit{paces}}
\def\Quotfunctor{\mathrm{Quot}}
\def\Hilbfunctor{\mathrm{Hilb}}
\def\Curvesstack{\mathcal{C}\!\mathit{urves}}
\def\Polarizedstack{\mathcal{P}\!\mathit{olarized}}
\def\Complexesstack{\mathcal{C}\!\mathit{omplexes}}
% \Pic is the operator that assigns to X its picard group, usage \Pic(X)
% \Picardstack_{X/B} denotes the Picard stack of X over B
% \Picardfunctor_{X/B} denotes the Picard functor of X over B
\def\Pic{\mathop{\mathrm{Pic}}\nolimits}
\def\Picardstack{\mathcal{P}\!\mathit{ic}}
\def\Picardfunctor{\mathrm{Pic}}
\def\Deformationcategory{\mathcal{D}\!\mathit{ef}}
\usepackage{cleveref}
\crefname{originallemma}{original lemma}{original lemmas}
\Crefname{originallemma}{Original lemma}{Original lemmas}
\makeatletter
\newlabel{lemma-local-lift-map}{{08FN}{source}{}{sourcepointer-1}{}}
\newlabel{lemma-local-lift-map@cref}{{[originallemma][0][]08FN}{[1][0][]source}{}{}{}}
\makeatother
\title{Maps from a strictly perfect complex to an acyclic complex are locally homotopic to zero}
\author{The Stacks Project Authors}
\date{stacks-project:sites-cohomology.tex:local-snapshot-20261011; GNU FDL 1.2 or later}
\begin{document}
\maketitle
\noindent\textit{Editorial note.} All delivered mathematical text is the author's own, copied from the pinned Stacks Project chapter file sites-cohomology.tex; no statement or proof step was rewritten, reordered or removed.\par\medskip
\noindent\textit{Editorial note.} The proof is the author's complete argument: part (1) is reduced to part (2) as the author states, and part (2) is a full induction whose base case and induction step are both written out.\par\medskip
\noindent\textit{Editorial note.} The single labelled reference occurring inside the counted spans is carried in assembly.external\_references as a non-load-bearing navigation pointer to the same chapter, so no load-bearing mathematics is replaced by a stub.\par\medskip
\noindent\textit{Editorial note.} Cross-chapter restatement: the same label occurs in chapter cohomology for a ringed space, and where applicable in chapter more-algebra for a ring. Those statements are the special cases of this ringed-site statement, so per the lane-s4 rule only this more general occurrence is counted here; the special-case twins are to be left uncounted by the lanes holding those chapters.\par\medskip
\noindent\textit{Editorial note.} The preamble is the chapter's own retained minimal preamble (sources/stacks-sites-cohomology/retained-preamble.tex), a verbatim excerpt of the native preamble: the Xy-pic setup and the native theorem declarations with the mathematical macros.\par\medskip
\noindent\textit{Allowed standard prerequisites (editorial).}\begin{itemize}
\item Strictly perfect complex of O-modules on a ringed site. Same chapter, the definition of a strictly perfect complex as a bounded complex of direct summands of finite free modules. Supplies the hypothesis on E\textasciicircum{}bullet that makes the length induction and the split exact sequence available.
\item K-flat and K-injective complexes and the homotopy category K(O\_U). Same chapter, the homotopy category of complexes of O\_U-modules and its triangulated structure. Supplies the distinguished triangle attached to the split exact sequence of complexes, whose long exact Hom sequence is used in the induction step.
\item Coverings of an object of a site and locality of homotopy. Same chapter, the passage to the members of a covering U\_i -> U and of restriction to the localizations C/U\_i. Supplies the meaning of being homotopic to zero on the members of a covering, which is the conclusion of both parts.
\end{itemize}\medskip
\ifdefined\hypertarget\hypertarget{sourcepointer-1}{}\fi
\noindent\textit{Source reference.} Original lemma 08FN: sources/stacks-sites-cohomology/pinned-original-source.tex, source lines 11772--11772. Same-chapter navigation pointer invoked at proof line 11820 for the base case of the induction, where a map out of a shifted summand of a finite free module is lifted locally. The pointer only resolves that reference for the reader; the cited lemma is not reproduced inside the counted span, and no bibliography citation is involved.\par\medskip
\begin{lemma}
\label{lemma-local-homotopy}
Let $(\mathcal{C}, \mathcal{O})$ be a ringed site. Let $U$ be an object
of $\mathcal{C}$.
\begin{enumerate}
\item Let $\alpha : \mathcal{E}^\bullet \to \mathcal{F}^\bullet$
be a morphism of complexes of $\mathcal{O}_U$-modules
with $\mathcal{E}^\bullet$ strictly perfect and $\mathcal{F}^\bullet$
acyclic. Then there exists a covering $\{U_i \to U\}$ such that each
$\alpha|_{U_i}$ is homotopic to zero.
\item Let $\alpha : \mathcal{E}^\bullet \to \mathcal{F}^\bullet$
be a morphism of complexes of $\mathcal{O}_U$-modules
with $\mathcal{E}^\bullet$ strictly perfect, $\mathcal{E}^i = 0$
for $i < a$, and $H^i(\mathcal{F}^\bullet) = 0$ for $i \geq a$.
Then there exists a covering $\{U_i \to U\}$ such that each
$\alpha|_{U_i}$ is homotopic to zero.
\end{enumerate}
\end{lemma}

\begin{proof}
The first statement follows from the second, hence we only prove (2).
We will prove this by induction on the length of the complex
$\mathcal{E}^\bullet$. If $\mathcal{E}^\bullet \cong \mathcal{E}[-n]$
for some direct summand $\mathcal{E}$ of a finite free
$\mathcal{O}$-module and integer $n \geq a$, then the result follows from
Lemma \ref{lemma-local-lift-map} and the fact that
$\mathcal{F}^{n - 1} \to \Ker(\mathcal{F}^n \to \mathcal{F}^{n + 1})$
is surjective by the assumed vanishing of $H^n(\mathcal{F}^\bullet)$.
If $\mathcal{E}^i$ is zero except for $i \in [a, b]$, then we have a
split exact sequence of complexes
$$
0 \to \mathcal{E}^b[-b] \to \mathcal{E}^\bullet \to
\sigma_{\leq b - 1}\mathcal{E}^\bullet \to 0
$$
which determines a distinguished triangle in
$K(\mathcal{O}_U)$. Hence an exact sequence
$$
\Hom_{K(\mathcal{O}_U)}(
\sigma_{\leq b - 1}\mathcal{E}^\bullet, \mathcal{F}^\bullet)
\to
\Hom_{K(\mathcal{O}_U)}(\mathcal{E}^\bullet, \mathcal{F}^\bullet)
\to
\Hom_{K(\mathcal{O}_U)}(\mathcal{E}^b[-b], \mathcal{F}^\bullet)
$$
by the axioms of triangulated categories. The composition
$\mathcal{E}^b[-b] \to \mathcal{F}^\bullet$ is homotopic to zero
on the members of a covering of $U$ by the above,
whence we may assume our map comes from an element in the
left hand side of the displayed exact sequence above. This element
is zero on the members of a covering of $U$ by induction hypothesis.
\end{proof}

\end{document}
